%!TEX root = tfg.tex
\chapter{Declaración de uso de inteligencia artificial}

\begin{abstract}
En este anexo se declara el uso de herramientas de inteligencia artificial durante el
desarrollo de \textit{Via Inferni} y durante la elaboración de esta memoria. Su utilización
ha tenido un carácter asistencial y de apoyo, manteniendo siempre la revisión, validación y
responsabilidad final por parte de los autores del proyecto.
\end{abstract}

\section{Finalidad del uso de la IA}
El uso de herramientas de inteligencia artificial ha estado orientado principalmente a dos
ámbitos: la revisión y redacción de texto en la memoria, y la asistencia al diseño e
implementación del videojuego en Unity.

En el caso de la memoria, la IA se ha utilizado como apoyo para mejorar la claridad de la
redacción, reorganizar explicaciones, revisar coherencia entre apartados y transformar
ideas técnicas en texto académico más estructurado. No se ha empleado como sustituto del
trabajo de análisis, sino como herramienta de apoyo para expresar con mayor precisión las
decisiones tomadas durante el proyecto.

En el desarrollo del videojuego, la IA ha servido como asistencia técnica para comprender
mejor Unity y sus flujos de trabajo. Aunque los autores contaban con conocimientos previos
de C\#, no habían utilizado Unity antes de este proyecto. Por ello, la IA se empleó para
consultar guías de uso, resolver dudas sobre componentes del motor, comprender patrones
habituales de implementación y aprovechar funcionalidades propias del entorno de Unity.

\section{Apoyo al diseño e implementación}
La asistencia proporcionada por la IA fue especialmente relevante durante las primeras fases
del desarrollo, donde el principal obstáculo no era el lenguaje de programación, sino el
desconocimiento del motor y de su ecosistema. Unity introduce conceptos propios, como
escenas, prefabs, componentes, ciclo de vida de objetos, físicas, serialización en el
editor, gestión de entradas, cámaras e interfaz, que requerían una curva de aprendizaje
específica.

En este contexto, la IA ayudó a superar el umbral inicial de habilidad necesario para
trabajar de forma productiva con Unity. Permitió obtener explicaciones guiadas, comparar
alternativas de diseño, detectar enfoques más adecuados para determinados sistemas y
aprovechar características del motor que, de otro modo, habrían requerido más tiempo de
investigación.

Esta ayuda resultó importante para alcanzar los objetivos del TFG, ya que permitió dedicar
más esfuerzo al diseño del videojuego, a la integración de funcionalidades y a la
documentación del proceso, reduciendo el tiempo invertido en resolver dudas básicas sobre
el uso de la herramienta.

\section{Revisión y responsabilidad sobre el contenido generado}
Todos los elementos generados o sugeridos mediante inteligencia artificial fueron revisados
posteriormente por los autores. Esta revisión incluyó la comprobación de que el contenido
fuera correcto, coherente con el proyecto y adecuado al nivel de detalle requerido en la
memoria.

En el caso del código o de las decisiones técnicas, las propuestas asistidas por IA no se
incorporaron de forma automática. Antes de integrarlas se analizaron, adaptaron y probaron
dentro del contexto real del proyecto. Del mismo modo, los textos generados o reformulados
fueron revisados para asegurar que reflejaban fielmente el trabajo realizado y las
decisiones tomadas por el equipo.

Por tanto, la IA se considera una herramienta de apoyo al aprendizaje, la documentación y
la productividad, pero no una fuente autónoma de decisión. La autoría, validación y
responsabilidad final del proyecto corresponden íntegramente a los autores del TFG.
